
\documentclass{article}
\usepackage{colortbl}
\usepackage{makecell}
\usepackage{multirow}
\usepackage{supertabular}

\begin{document}

\newcounter{utterance}

\centering \large Interaction Transcript for game `cladder', experiment `full\_v1.5\_default', episode 3752 with selfrefine{-}llama{-}3.1{-}8b{-}instruct.
\vspace{24pt}

{ \footnotesize  \setcounter{utterance}{1}
\setlength{\tabcolsep}{0pt}
\begin{supertabular}{c@{$\;$}|p{.15\linewidth}@{}p{.15\linewidth}p{.15\linewidth}p{.15\linewidth}p{.15\linewidth}p{.15\linewidth}}
    \# & \multicolumn{2}{c}{Player} && \multicolumn{2}{c}{Game Master} \\
    \hline

    \theutterance \stepcounter{utterance}  
    & & & \multicolumn{4}{p{0.6\linewidth}}{
        \cellcolor[rgb]{0.9,0.9,0.9}{
            \makecell[{{p{\linewidth}}}]{
                \texttt{\tiny{[P1$\langle$GM]}}
                \texttt{You are an expert in causal inference. The following question is not a typical commonsense query, but rather a meticulously designed question created by a professor specializing in causal inference, intended to assess the students' mastery of the course content.} \\
\texttt{Imagine a self{-}contained, hypothetical world with only the following conditions, and without any unmentioned factors or causal relationships: Listening to jazz has a direct effect on elite institution admission status. Effort has a direct effect on elite institution admission status. For students accepted to elite institutions, the correlation between listening to jazz and being hard{-}working is {-}0.09. If we look at students accepted to elite institutions, does it mean that listening to jazz affects being hard{-}working?} \\
\texttt{Start your answer with 'Yes' or 'No', followed by additional reasoning or evidence to support your explanation.} \\
            }
        }
    }
     \\ \\

    \theutterance \stepcounter{utterance}  
    & & \multicolumn{4}{p{0.6\linewidth}}{
        \cellcolor[rgb]{0.95,0.95,0.95}{
            \makecell[{{p{\linewidth}}}]{
                \texttt{\tiny{[P1$\leftrightarrow$P]}}
                \texttt{{[}DRAFT{]}} \\
\texttt{No.} \\
\\ 
\texttt{The correlation between listening to jazz and being hard{-}working among students accepted to elite institutions is {-}0.09, indicating a negative relationship. However, this does not imply causation from listening to jazz to being hard{-}working. } \\
\\ 
\texttt{In this hypothetical world, we have two direct effects on elite institution admission status: listening to jazz and effort. The negative correlation between listening to jazz and being hard{-}working among students accepted to elite institutions can be explained by the fact that both variables are correlated with a third variable, elite institution admission status, but through different causal pathways.} \\
\\ 
\texttt{To illustrate this, let's consider a possible scenario: students who are naturally more hard{-}working tend to be less interested in jazz, while students who are naturally more interested in jazz tend to be less hard{-}working. However, both hard{-}working students and jazz{-}listening students are more likely to be accepted to elite institutions due to their respective characteristics.} \\
\\ 
\texttt{In this case, the negative correlation between listening to jazz and being hard{-}working among students accepted to elite institutions is a result of confounding by the underlying causal relationships, rather than a causal effect of listening to jazz on being hard{-}working. Therefore, we cannot conclude that listening to jazz affects being hard{-}working based on this correlation.} \\
            }
        }
    }
    & \\ \\

    \theutterance \stepcounter{utterance}  
    & & \multicolumn{4}{p{0.6\linewidth}}{
        \cellcolor[rgb]{0.95,0.95,0.95}{
            \makecell[{{p{\linewidth}}}]{
                \texttt{\tiny{[P1$\leftrightarrow$P]}}
                \texttt{{[}SELF{-}FEEDBACK{]}} \\
\texttt{FORMAT: ok} \\
\texttt{RULES: ok} \\
\texttt{STRATEGY: ok} \\
\texttt{REASON: The answer correctly identifies the negative correlation as a result of confounding by the underlying causal relationships, rather than a causal effect of listening to jazz on being hard{-}working.} \\
            }
        }
    }
    & \\ \\

    \theutterance \stepcounter{utterance}  
    & \multicolumn{4}{p{0.6\linewidth}}{
        \cellcolor[rgb]{0.561,0.737,0.561}{
            \makecell[{{p{\linewidth}}}]{
                \texttt{\tiny{[P1$\rangle$GM]}}
                \texttt{No.} \\
\\ 
\texttt{The correlation between listening to jazz and being hard{-}working among students accepted to elite institutions is {-}0.09, indicating a negative relationship. However, this does not imply causation from listening to jazz to being hard{-}working. } \\
\\ 
\texttt{In this hypothetical world, we have two direct effects on elite institution admission status: listening to jazz and effort. The negative correlation between listening to jazz and being hard{-}working among students accepted to elite institutions can be explained by the fact that both variables are correlated with a third variable, elite institution admission status, but through different causal pathways.} \\
\\ 
\texttt{To illustrate this, let's consider a possible scenario: students who are naturally more hard{-}working tend to be less interested in jazz, while students who are naturally more interested in jazz tend to be less hard{-}working. However, both hard{-}working students and jazz{-}listening students are more likely to be accepted to elite institutions due to their respective characteristics.} \\
\\ 
\texttt{In this case, the negative correlation between listening to jazz and being hard{-}working among students accepted to elite institutions is a result of confounding by the underlying causal relationships, rather than a causal effect of listening to jazz on being hard{-}working. Therefore, we cannot conclude that listening to jazz affects being hard{-}working based on this correlation.} \\
            }
        }
    }
    & & \\ \\

    \theutterance \stepcounter{utterance}  
    & & & \multicolumn{2}{p{0.3\linewidth}}{
        \cellcolor[rgb]{0.9,0.9,0.9}{
            \makecell[{{p{\linewidth}}}]{
                \texttt{\tiny{[GM$|$GM]}}
                \texttt{No} \\
            }
        }
    }
    & & \\ \\

    \theutterance \stepcounter{utterance}  
    & & & \multicolumn{2}{p{0.3\linewidth}}{
        \cellcolor[rgb]{0.9,0.9,0.9}{
            \makecell[{{p{\linewidth}}}]{
                \texttt{\tiny{[GM$|$GM]}}
                \texttt{no} \\
            }
        }
    }
    & & \\ \\

    \theutterance \stepcounter{utterance}  
    & & & \multicolumn{2}{p{0.3\linewidth}}{
        \cellcolor[rgb]{0.9,0.9,0.9}{
            \makecell[{{p{\linewidth}}}]{
                \texttt{\tiny{[GM$|$GM]}}
                \texttt{game\_result = WIN} \\
            }
        }
    }
    & & \\ \\

\end{supertabular}
}

\end{document}
